\label{sec:legal}

\subsection{Section~203 and Section~2: Distinct but Overlapping Frameworks}

Section~203 and Section~2 of the Voting Rights Act protect overlapping
but distinct interests. Section~2 protects minority voting strength:
it prohibits redistricting plans that dilute minority communities'
ability to elect representatives of their choice, and through the
\textit{Thornburg v.\ Gingles}\footnote{\textit{Thornburg v.\ Gingles},
478 U.S.\ 30 (1986).} framework, it requires majority-minority districts
in jurisdictions where the three Gingles preconditions are met.
Section~203 protects access: it ensures that language minority citizens
can receive election materials in their native language, enabling effective
participation in the electoral process.

A jurisdiction can be covered by both provisions---as many are. A
Spanish-speaking community in South Texas may be both a Section~2
majority-Hispanic community (requiring a majority-Hispanic district)
and a Section~203-covered community (requiring Spanish-language
election materials). These requirements reinforce each other, but
they also interact in ways that redistricting must address.

\subsection{Jurisdictions with Both Section~203 and Section~2 Coverage}

This paper identifies 23 counties nationally that, as of 2023, have both
Section~203 coverage for a language minority group and a Gingles-feasible
majority-minority community within the same language group. These counties
present the most demanding redistricting challenge: the district plan must
simultaneously satisfy Section~2's majority-minority district requirement
and avoid splitting the Section~203-covered community between districts.

The counties divide into three categories:

\textbf{Urban Spanish-speaking counties} (Bexar, El Paso, Webb, Hidalgo
in Texas; Bernalillo in New Mexico; Miami-Dade in Florida): Large,
geographically concentrated Hispanic communities with both Section~203
Spanish coverage and Section~2 majority-Hispanic district feasibility.
These are the least problematic: the community's geographic concentration
makes it natural for both Section~2 and Section~203 purposes to keep it
within one district.

\textbf{Multi-language urban counties} (Santa Clara, Los Angeles, Alameda
in California; King County in Washington; Maricopa in Arizona): Counties
with Section~203 coverage in multiple languages, at least one of which
also has Section~2 feasibility. The challenge here is that different
language communities may be geographically separated within the county,
making it difficult to satisfy Section~2 for one community without
splitting another community covered by Section~203.

\textbf{Native American reservation counties} (Apache, Navajo in Arizona;
San Juan in New Mexico; various in Alaska): Tribal nations whose members
are covered by Section~203 for Native American or Alaska Native languages
and in many cases are also majority populations with Section~2 feasibility.
The geographic structure of reservations---often non-contiguous with state
borders---presents unique redistricting challenges.

\subsection{How \textsc{bisect} VRA Mode Handles Both Requirements}

The \textsc{bisect} VRA mode (\texttt{--weights-override vra-aligned}) is
designed to satisfy Section~2's majority-minority district requirement by
assigning higher METIS edge weights to boundaries between tracts with high
minority population concentration, making the algorithm reluctant to split
minority communities. This directly serves Section~203 community integrity
as well: keeping minority tracts together serves both VRA provisions.

The interaction is not always perfectly aligned. In a multi-language county,
the VRA mode optimizes for the language group that meets the Section~2
Gingles threshold---typically the largest minority group---and may not
specifically optimize for a second language minority community that has
Section~203 coverage but not Section~2 feasibility. In these cases, the
\textsc{bisect} county-weight mode (\texttt{--weights-override county})
provides better Section~203 protection than the VRA mode, because county
integrity protects the administrative unit at which Section~203 coverage
is determined.

\subsection{Documentation for Court Filings}

Expert witnesses presenting \textsc{bisect} plans in jurisdictions with
both Section~203 and Section~2 requirements should document:

\begin{enumerate}
\item \textbf{Section~203 coverage determination}: Cite the applicable
  Federal Register determination (88 Fed.\ Reg.\ 34,396 for the 2023
  determination) and identify which counties in the state are covered
  and for which languages.

\item \textbf{County split comparison}: Present Table~\ref{tab:results-203}
  or equivalent data showing that the \textsc{bisect} plan produces fewer
  Section~203 county splits than the enacted plan. This demonstrates
  that the algorithmic plan is more administratively compatible with
  Section~203 compliance.

\item \textbf{Community integrity scores}: For each language minority
  community with Section~203 coverage, report the community integrity
  score for both the \textsc{bisect} plan and the enacted plan. A
  significantly higher score for the \textsc{bisect} plan is evidence
  that the enacted plan's splits are not compelled by population arithmetic.

\item \textbf{Section~2 compliance}: Confirm that the \textsc{bisect}
  plan satisfies Section~2's majority-minority district requirements
  in all applicable jurisdictions, and that the Section~203-motivated
  design choices (county integrity, community compactness) are compatible
  with the Section~2 requirements.
\end{enumerate}

\subsection{The Gingles--Section~203 Interaction: A Note on Legal Hierarchy}

When Section~203 community integrity and Section~2 majority-minority
district requirements conflict---as they may in multi-language counties
where achieving a majority-minority district for one group splits a
Section~203-covered community of another---Section~2's requirements are
legally paramount. The Supreme Court has been clear that majority-minority
districts must be drawn where the Gingles preconditions are met.\footnote{
See \textit{Alabama Legislative Black Caucus v.\ Alabama}, 575 U.S.\ 254 (2015)
(holding that racial predominance in redistricting requires strict scrutiny
but that VRA-required majority-minority districts satisfy that scrutiny).}

Section~203, by contrast, creates an administrative obligation on the
county election office, not a constraint on district lines. A district plan
that splits a Section~203-covered community does not violate Section~203
as long as the county election office continues to provide language
assistance county-wide. The Section~203 consideration is therefore
a secondary optimization criterion: after Section~2 requirements are satisfied,
the plan should minimize Section~203 administrative burden by minimizing
covered county splits and maximizing community integrity.

This hierarchy is reflected in the \textsc{bisect} pipeline: the VRA mode
(Section~2) takes precedence, and within the VRA-constrained solution space,
the county-weight mode (Section~203) further optimizes for county integrity.
